% ===================================================================
%  TABLEAU N° 9 — La Montagne. — La Ville d'eaux. / La Forêt. — La Chasse.
%  Feuillets 56 a 61 du fac-simile = folios 54 a 59.
%  Correspondance : folio imprime = numero de feuillet - 2.
%
%  Contrairement au tableau io correspondant (texto/io/18-tabelo-09.tex),
%  le francais ne porte AUCUN des intertitres decoratifs qui separent
%  les alineas en io ("Monto-peizajo.", "L'Exkursionero.", etc.) : seuls
%  les intitules de scene ("Premiere scene." / "Deuxieme scene." et leurs
%  titres) sont communs aux deux editions. Les cles %%K suivent donc une
%  numerotation continue d'alineas a l'interieur de chaque scene.
%
%  Quelques renvois numeriques entre parentheses divergent entre les
%  deux editions sur ce tableau (ex. « bohémien (11) » cote francais,
%  contre « cigano (41) » cote ido ; « pâturage (19) » contre
%  « pastureyo (49) »). Ce sont des ecarts REELS du fac-simile,
%  releves tels quels, non des erreurs de transcription.
% ===================================================================

% --- Feuillet 56 (folio 54) -----------------------------------------
\begin{VUpage}[56]{54}
%%K t09-apar-1 apar
\VUsaut{4.0mm}
\VUtitre{15.0pt}{60}{TABLEAU N\textsuperscript{o} 9}
\VUsaut{3.0mm}

%%K t09-tit-1 sub
\VUcentre{12.0pt}{0}{\textit{Première scène.}}
\VUsaut{3.0mm}

%%K t09-tit-2 sub
\VUpk{11.4pt}{40}{La Montagne. --- La Ville d'eaux.}
\VUsaut{3.0mm}

%%K t09-c1-01-1 p
\VUblancAlinea
1. --- Depuis huit jours je suis à Pratz. Je ne puis\nl
exprimer toute l'admiration que j'ai éprouvée lors\cc
que je me suis trouvé en face de la merveilleuse\nl
\VUgras{chaîne} \textsuperscript{(1)} des Pyrénées dont la \VUgras{crête} \textsuperscript{(2)} se découpe\nl
sur le fond bleu du ciel comme un gigantesque feston.

%%K t09-c1-02-1 p
\VUblancAlinea
2. --- Je ne m'imaginais pas que les \VUgras{sommets} \textsuperscript{(3)}\nl
pyrénéens pouvaient atteindre une \VUgras{altitude} aussi\nl
grande, et je suis resté saisi d'étonnement devant les\nl
magnifiques \VUgras{glaciers} \textsuperscript{(5)} et les \VUgras{pics} \textsuperscript{(4)} couverts de\nl
neige sur lesquels roulent de si terribles \VUgras{avalanches.}

%%K t09-c1-03-1 p
\VUblancAlinea
3. --- Dès le lendemain de mon arrivée je suis allé\nl
au vieux \VUgras{volcan} \textsuperscript{(6)} éteint qui se trouve près de Pratz.\nl
Sur les bords du \VUgras{cratère,} arides et dénudés, on voit\nl
encore parfaitement les traces du passage de la \VUgras{lave}\nl
qui coulait, il y a plusieurs siècles, sur le \VUgras{flanc de la}\nl
\VUgras{montagne} \textsuperscript{(7)}. J'ai pu retrouver, sur l'un des \VUgras{ver}\cc
\VUgras{sants,} quelques parcelles de cette lave.

%%K t09-c1-04-1 p
\VUblancAlinea
4. --- Pratz est situé sur la frontière, et la \VUgras{forte}\cc
\VUgras{resse} \textsuperscript{(8)} qui défend le \VUgras{col} \textsuperscript{(27)} séparant la France de\nl
l'Espagne rappelle tout à fait ces vieux châteaux des\nl
bords du Rhin qui semblent guetter une \VUgras{proie} du\nl
haut de leur rocher.

%%K t09-c1-05-1 p
\VUblancAlinea
5. --- Un \VUgras{aigle} \textsuperscript{(9)} énorme, qui a son aire non loin du\nl
\VUgras{fort} \textsuperscript{(8)}, attire souvent mon attention. Cet \VUgras{oiseau de}\nl
\VUgras{proie,} au \VUgras{bec} puissant, rapporte fréquemment à ses\nl
petits un agneau ou un chevreau qu'il tient dans ses\nl
\VUgras{serres.} La grandeur de ses \VUgras{ailes} est telle qu'il peut\nl
planer assez longtemps avec son lourd fardeau.

%%K t09-c1-06-1 p
\VUblancAlinea
6. --- La promenade la plus agréable à faire ici est\nl
l'excursion à la \VUgras{cascade} \textsuperscript{(10)} de Cau. La \VUgras{chute d'eau}\nl
tombe en bouillonnant et vient se perdre en un joli
\parplein
\end{VUpage}

% --- Feuillet 57 (folio 55) -----------------------------------------
\begin{VUpage}[57]{55}
%%K t09-c1-06-1 p suite
\VUcontinue
ruisseau dans le \VUgras{bois de sapin} \textsuperscript{(11)} situé à l'ouest de\nl
la ville. Nous sommes montés à travers ce bois à\nl
l'\VUgras{observatoire météorologique} \textsuperscript{(12)} et du haut du\nl
\VUgras{belvédère,} nous avons embrassé le \VUgras{panorama}\nl
entier de la région. Le \VUgras{paysage} est superbe et la\nl
\VUgras{nature} semble avoir prodigué au pays tous les\nl
trésors dont elle dispose.

%%K t09-c1-07-1 p
\VUblancAlinea
7. --- Pour redescendre, nous avons pris le \VUgras{funicu}\cc
\VUgras{laire} \textsuperscript{(13)} et nous nous sommes arrêtés à une petite\nl
\VUgras{station} près des \VUgras{ruines} \textsuperscript{(14)} d'un vieux \VUgras{château}\nl
\VUgras{féodal} habité autrefois par les seigneurs de Pratz.

%%K t09-c1-08-1 p
\VUblancAlinea
8. --- Pauvre château ! Que doit-il penser en voyant\nl
à ses pieds la coquette \VUgras{ville d'eaux} \textsuperscript{(15)} qui est main\cc
tenant une \VUgras{station thermale} à la mode ?

%%K t09-c1-08-2 p
\VUblancAlinea
Le \VUgras{climat} est si doux, l'\VUgras{eau minérale} si bienfai\cc
sante, que la \VUgras{cure} que l'on vient faire au \VUgras{sanato}\cc
\VUgras{rium} \textsuperscript{(17)} est un véritable plaisir.

%%K t09-c1-09-1 p
\VUblancAlinea
9. --- D'ailleurs, on a tout prodigué pour rendre le\nl
séjour agréable. Un grand \VUgras{viaduc} \textsuperscript{(16)} a été construit\nl
pour permettre d'établir une voie ferrée ; le \VUgras{parc} \textsuperscript{(18)}\nl
de l'\VUgras{établissement thermal} dans lequel se donnent\nl
les \VUgras{concerts} est aménagé d'une façon charmante.\nl
De jolies \VUgras{villas} \textsuperscript{(19)} ont été bâties autour du \VUgras{casino} \textsuperscript{(21)}\nl
et une \VUgras{chaussée} \textsuperscript{(22)} très bien entretenue rend les\nl
communications extrêmement faciles.

%%K t09-c1-10-1 p
\VUblancAlinea
10. --- Un simple \VUgras{talus} \textsuperscript{(23)} sépare la \VUgras{route} \textsuperscript{(20)} de la\nl
\VUgras{vallée} \textsuperscript{(25)} proprement dite, au fond de laquelle s'étale\nl
un joli petit \VUgras{lac} \textsuperscript{(26)} qui alimente un \VUgras{ruisseau} \textsuperscript{(24)} très\nl
clair.

%%K t09-c1-11-1 p
\VUblancAlinea
11. --- De l'autre côté du vallon, une \VUgras{mine} \textsuperscript{(28)} de\nl
fer est en exploitation. Je suis allé plusieurs fois voir\nl
les \VUgras{mineurs} descendre dans le \VUgras{puits,} et j'ai même\nl
pénétré dans la \VUgras{galerie} où ils passent leur journée\nl
à extraire le \VUgras{minerai} qui contient le \VUgras{métal.}

%%K t09-c1-12-1 p
\VUblancAlinea
12. --- Une importante \VUgras{fonderie} a été bâtie près de\nl
la mine pour utiliser immédiatement les richesses qui
\parplein
\end{VUpage}

% --- Feuillet 58 (folio 56) -----------------------------------------
\begin{VUpage}[58]{56}
%%K t09-c1-12-1 p suite
\VUcontinue
en sont retirées. Le \VUgras{haut fourneau} \textsuperscript{(29)}, chauffé par\nl
le \VUgras{charbon de terre}, qui abonde ici, permet d'obtenir\nl
la \VUgras{fonte} d'une façon rapide et peu coûteuse.

%%K t09-c1-13-1 p
\VUblancAlinea
13. --- Non loin de la mine on creuse une \VUgras{car}\cc
\VUgras{rière} \textsuperscript{(30)}. Ce n'est ni du \VUgras{granit}, ni du \VUgras{porphyre}, ni\nl
même du \VUgras{grès} que le vieux \VUgras{carrier} détache avec\nl
son \VUgras{pic}, mais de la simple \VUgras{pierre de taille} pour la\nl
construction des maisons.

%%K t09-c1-14-1 p
\VUblancAlinea
14. --- Un des plus jolis ornements de ce pays est\nl
certainement le \VUgras{sapin} \textsuperscript{(31)}. Partout, dans le fond d'un\nl
\VUgras{ravin} \textsuperscript{(32)}, au bord d'un \VUgras{torrent} \textsuperscript{(33)}, d'un \VUgras{gouffre} \textsuperscript{(34)}\nl
ou d'un précipice, ses fines aiguilles mettent leur\nl
note verte et gaie.

%%K t09-c1-15-1 p
\VUblancAlinea
15. --- Je profite de mes vacances pour faire de lon\cc
gues \VUgras{promenades.} Les \VUgras{chemins} \textsuperscript{(35)} qui serpentent\nl
dans la montagne permettent de parcourir, sans\nl
s'apercevoir de la distance, plusieurs \VUgras{kilomètres}\nl
et souvent plusieurs \VUgras{lieues.}

%%K t09-c1-15-2 p
\VUblancAlinea
Des \VUgras{bornes} \textsuperscript{(36)} et des \VUgras{poteaux indicateurs} \textsuperscript{(37)}\nl
jalonnent les routes, sur lesquelles on rencontre sou\cc
vent un jeune \VUgras{montagnard} \textsuperscript{(38)}, la \VUgras{besace} \textsuperscript{(40)} au\nl
côté, portant une \VUgras{marmotte} \textsuperscript{(39)}, ou bien quelque\nl
\VUgras{bohémien} \textsuperscript{(11)}, dont la \VUgras{toque de fourrure} \textsuperscript{(42)} con\cc
traste singulièrement avec le vieux \VUgras{capuchon} \textsuperscript{(43)}\nl
déchiré. Il conduit par sa \VUgras{chaîne} un \VUgras{ours} \textsuperscript{(14)} savant.

%%K t09-c1-16-1 p
\VUblancAlinea
16. --- Presque tous les jours je vais avec un\nl
\VUgras{guide} \textsuperscript{(48)} faire une \VUgras{ascension}, et je suis très fier\nl
lorsque, le \VUgras{sac} \textsuperscript{(47)} au dos, l'\VUgras{alpenstock} à la main,\nl
en vrai \VUgras{touriste} \textsuperscript{(46)}, j'escalade les \VUgras{rochers} \textsuperscript{(45)}.

%%K t09-c1-17-1 p
\VUblancAlinea
17. --- Du haut d'une montagne, les habitants de la\nl
plaine ne semblent pas plus gros que des fourmis ; un\nl
\VUgras{troupeau de moutons} \textsuperscript{(50)} paissant dans un \VUgras{pâtu}\cc
\VUgras{rage} \textsuperscript{(19)} paraît un jouet d'enfant. Un \VUgras{pin} \textsuperscript{(51)} ou même\nl
un \VUgras{chalet} \textsuperscript{(52)} font à peine une tache dans la verdure.

%%K t09-c1-18-1 p
\VUblancAlinea
18. --- Pendant ces excursions, nous croisons fré\cc
quemment dans le \VUgras{sentier} \textsuperscript{(53)} un \VUgras{chasseur de}
\parplein
\end{VUpage}

% --- Feuillet 59 (folio 57) -----------------------------------------
\begin{VUpage}[59]{57}
%%K t09-c1-18-1 p suite
\VUcontinue
\VUgras{chamois} \textsuperscript{(54)} rapportant sur son épaule l'animal qu'il\nl
vient de tuer.

%%K t09-c1-19-1 p
\VUblancAlinea
19. --- J'ai mis dans mon herbier une branche de\nl
\VUgras{fougère} \textsuperscript{(55)} très curieuse et un brin de \VUgras{bruyère} \textsuperscript{(57)}\nl
rose fort joli, que j'ai cueillis à l'entrée d'une petite\nl
\VUgras{grotte} \textsuperscript{(56)}, dans ma dernière promenade.

%%K t09-tit-3 sub
\VUsaut{3.0mm}
\VUcentre{12.0pt}{0}{\textit{Deuxième scène.}}
\VUsaut{3.0mm}

%%K t09-tit-4 sub
\VUpk{11.4pt}{40}{La Forêt. --- La Chasse.}
\VUsaut{3.0mm}

%%K t09-c2-01-1 p
\VUblancAlinea
1. --- J'ai passé hier une délicieuse journée dans la\nl
forêt. L'activité qui règne parmi les animaux et les\nl
\VUgras{insectes} \textsuperscript{(5)} dont elle est peuplée m'a vivement\nl
intéressé.

%%K t09-c2-01-2 p
\VUblancAlinea
L'\VUgras{araignée} \textsuperscript{(1)} qui tisse sa \VUgras{toile} \textsuperscript{(2)} entre deux\nl
branches pour attraper la \VUgras{mouche} \textsuperscript{(3)} qui passe,\nl
l'\VUgras{escargot} \textsuperscript{(4)} portant péniblement sa coquille, le\nl
\VUgras{cerf-volant} à la recherche de sa proie, la \VUgras{fourmi}\nl
diligente s'agitant auprès de la \VUgras{fourmilière} \textsuperscript{(6)}, tout\nl
ce petit monde allait, venait, travaillait avec une\nl
animation extraordinaire.

%%K t09-c2-02-1 p
\VUblancAlinea
2. --- Un \VUgras{serpent} \textsuperscript{(7)}, que j'avais pris d'abord pour\nl
une \VUgras{vipère,} mais qui n'était qu'une simple \VUgras{cou}\cc
\VUgras{leuvre,} était blotti sous un tronc d'arbre, et sortait de\nl
temps en temps sa petite tête fine qui semblait toujours\nl
prête à mordre. Devant moi, une grosse \VUgras{limace} \textsuperscript{(8)}\nl
rampait lourdement après avoir dévoré une des déli\cc
cieuses petites \VUgras{fraises} \textsuperscript{(11)} qui croissent ici en abon\cc
dance.

%%K t09-c2-03-1 p
\VUblancAlinea
3. --- Le sol était parsemé de \VUgras{fleurs des bois} ; des\nl
\VUgras{primevères} \textsuperscript{(9)}, des \VUgras{pervenches} \textsuperscript{(10)} et un grand\nl
nombre d'autres plantes que je ne connaissais pas,\nl
fleurissaient au milieu des \VUgras{touffes d'herbe} \textsuperscript{(12)}.

%%K t09-c2-04-1 p
\VUblancAlinea
4. --- Dans cette région, la \VUgras{truffe} est presque aussi\nl
commune que le \VUgras{champignon} \textsuperscript{(14)}, et un petit \VUgras{porc} \textsuperscript{(13)}\nl
appartenant à un forestier cherchait d'un air gour\cc
mand s'il n'en découvrirait pas quelques-unes.
\end{VUpage}

% --- Feuillet 60 (folio 58) -----------------------------------------
\begin{VUpage}[60]{58}
%%K t09-c2-05-1 p
\VUblancAlinea
5. --- J'ai assisté à la fin d'une scène qui m'a bien\nl
amusé. Grâce au flair de son \VUgras{chien basset} \textsuperscript{(15)}, un\nl
\VUgras{garde forestier} \textsuperscript{(16)} venait de surprendre un \VUgras{bra}\cc
\VUgras{connier} \textsuperscript{(22)} au moment où celui-ci se préparait à\nl
emporter le \VUgras{gibier} \textsuperscript{(17)} qu'il avait tué.

%%K t09-c2-05-2 p
\VUblancAlinea
Préférant abandonner son butin plutôt que de se\nl
laisser arrêter, le braconnier s'était enfui, et \VUgras{liè}\cc
\VUgras{vre} \textsuperscript{(18)}, \VUgras{biche} \textsuperscript{(19)}, \VUgras{chevreuil} \textsuperscript{(20)} et \VUgras{sanglier} \textsuperscript{(21)}\nl
gisaient sur l'herbe à la grande satisfaction du garde.

%%K t09-c2-06-1 p
\VUblancAlinea
6. --- La variété d'arbres que renferme la forêt est\nl
étonnante. Les \VUgras{hêtres} \textsuperscript{(23)} et les \VUgras{bouleaux} \textsuperscript{(26)}, les\nl
\VUgras{pins,} qui produisent la \VUgras{résine} \textsuperscript{(27)}, les \VUgras{chênes} \textsuperscript{(29)} et\nl
beaucoup d'autres encore entremêlent leurs branches.

%%K t09-c2-06-2 p
\VUblancAlinea
Le \VUgras{lierre} \textsuperscript{(24)}, qui s'attache au tronc des grands\nl
arbres, donne à la \VUgras{futaie} \textsuperscript{(25)} une teinte verte brillante,\nl
qui s'unit agréablement au ton plus clair et plus\nl
doux de la \VUgras{mousse} \textsuperscript{(31)} dans laquelle le \VUgras{pic-vert} \textsuperscript{(30)}\nl
cherche des insectes.

%%K t09-c2-07-1 p
\VUblancAlinea
7. --- Les vieux chênes m'ont charmé. Leurs grosses\nl
\VUgras{racines} \textsuperscript{(32)} tordues, sortant à moitié de terre, leur\nl
donne une apparence magnifique de force et de\nl
vigueur, et je me demande comment le \VUgras{proprié}\cc
\VUgras{taire} \textsuperscript{(35)} peut se résoudre à les faire abattre. et à les\nl
livrer à la \VUgras{scie} \textsuperscript{(34)} du \VUgras{bûcheron} \textsuperscript{(33)}. Les pauvres\nl
\VUgras{troncs} \textsuperscript{(36)} dépouillés de leur \VUgras{écorce} \textsuperscript{(37)}, ou fendus\nl
par la \VUgras{cognée} \textsuperscript{(38)} du \VUgras{journalier} \textsuperscript{(39)}, me font peine à\nl
voir.

%%K t09-c2-08-1 p
\VUblancAlinea
8. --- Comme me l'expliquait un vieux \VUgras{charbon}\cc
\VUgras{nier} \textsuperscript{(41)} qui préparait avec sa \VUgras{pelle} une \VUgras{meule} \textsuperscript{(42)}\nl
de \VUgras{charbon}, cela permet aux indigents de compléter\nl
leur \VUgras{fagot} \textsuperscript{(40)} de \VUgras{bois mort}, avec les \VUgras{brindilles} des\nl
arbres abattus, et personne ne s'en plaint.

%%K t09-c2-09-1 p
\VUblancAlinea
9. --- Je voulais aller me reposer un moment à la\nl
\VUgras{maison forestière} \textsuperscript{(46)}, mais lorsque je suis arrivé\nl
auprès de l'\VUgras{étang} \textsuperscript{(43)}, sur lequel nage un beau\nl
\VUgras{cygne} \textsuperscript{(45)}, je me suis aperçu qu'une récente \VUgras{inon}\cc
\VUgras{dation} \textsuperscript{(44)} avait converti le chemin en un véritable
\parplein
\end{VUpage}

% --- Feuillet 61 (folio 59) -----------------------------------------
\begin{VUpage}[61]{59}
%%K t09-c2-09-1 p suite
\VUcontinue
\VUgras{marécage.} Je dus faire un détour et, en passant\nl
près du \VUgras{cèdre} \textsuperscript{(49)}, des \VUgras{peupliers} \textsuperscript{(48)} et de l'\VUgras{orme} \textsuperscript{(47)}\nl
qui sont sur la lisière de la forêt, je vis déboucher\nl
d'un \VUgras{taillis} \textsuperscript{(54)} un magnifique \VUgras{cerf} \textsuperscript{(50)} poursuivi par\nl
une \VUgras{meute} \textsuperscript{(51)} acharnée qu'excitait encore le \VUgras{cor} \textsuperscript{(53)}\nl
des \VUgras{piqueurs} \textsuperscript{(52)}. La pauvre bête était à bout de\nl
forces. Elle vint tomber mourante à quelques pas de\nl
moi.

%%K t09-c2-10-1 p
\VUblancAlinea
10. --- Le fils du garde forestier, que le bruit de la\nl
\VUgras{chasse à courre} avait attiré, sortit de la maison\nl
et nous retournâmes tous deux dans la forêt, où il\nl
avait aperçu sur la \VUgras{branche} d'un vieux chêne une\nl
\VUgras{pie} \textsuperscript{(56)}. Il pensait que son nid devait être caché dans\nl
le \VUgras{feuillage} \textsuperscript{(58)} et il voulait le dénicher.

%%K t09-c2-10-2 p
\VUblancAlinea
Leste comme un \VUgras{écureuil} \textsuperscript{(61)}, il grimpa de \VUgras{bran}\cc
\VUgras{che} \textsuperscript{(55)} en branche, mais il ne trouva rien et ne réussit\nl
qu'à faire envoler une vieille \VUgras{corneille} \textsuperscript{(60)} perchée\nl
sur un \VUgras{rameau} \textsuperscript{(57)}.

\VUsaut{4.0mm}
\VUfilet{22mm}
\end{VUpage}
